%% sections/b-section-by-section.tex - H. R. 9510 Bill v4.0 (full bill) - Appendix B
%% Genre: per-section annotation/gloss. Source deliverable 02. Plain-English gloss.
\billsec{APPENDIX B.}{SECTION-BY-SECTION ANALYSIS (NON-OPERATIVE).}\phantomsection\label{toc:aB}

\uscprov{1}{This appendix walks through H. R. 9510 provision by provision and gives a
plain-English gloss of each part. It is an annotation, not operative text: where the
gloss and the enacted words differ, the words of section 3 and the conforming
amendments govern. The bill amends the Federal Food, Drug, and Cosmetic Act (21
U.S.C. 301 et seq.), current through Public Law 119-93, and its operative mechanism is
a single sequencing rule in a new section 515D (21 U.S.C. 360e-5), inserted after
section 515C (21 U.S.C. 360e-4): an automated verification, validation, and
uncertainty quantification process must clear robot-patient interaction code before
that code is generated or executed in a Physical AI oncology clinical investigation
\cite{usc-fdca, pl-119-93, usc-360e4}.}

\uscprov{1}{\textbf{Sec. 1. Short Title; Table of Contents.} Subsection (a) gives the
citation, the ``Verification Before Generation in Physical AI Oncology Trials Act of
2026.'' Subsection (b) lays out the table of contents and maps what follows. This is
standard front matter that fixes how the Act is cited and carries no substantive
duty.}

\uscprov{1}{\textbf{Sec. 2. Findings.} Section 2 records fourteen congressional
findings that justify the rule and show it is feasible in 2026. They state the problem
(robots in direct physical contact with patients, behavior at contact governed by
increasingly machine-generated software, no Federal rule that the code be verified
before it is generated or executed); explain why embodiment raises the stakes (a
disembodied error yields a wrong number a human can catch downstream, while an
embodied error yields a wrong motion no human can catch in time); identify the closest
existing analogues (the section 515C change control plan authority and FDA guidance,
the Quality Management System Regulation at 21 CFR part 820 in force since February 2,
2026, and 21 CFR part 11), noting each is voluntary and not required before generation;
and ground the mechanism in consensus standards, in human-over-AI precedent in State
law and the Medicare Advantage guardrail, and in Federal nondiscrimination law. The
findings supply the predicate a court or agency would consult when construing the Act
\cite{usc-360e4, fr-qmsr-2024, cfr-part11}.}

\uscprov{1}{\textbf{Sec. 3. Amendment to the Federal Food, Drug, and Cosmetic Act.}
The operative heart of the bill. Subsection (a) inserts the new section 515D into
subchapter V after section 515C; subsection (b) makes the ten conforming amendments;
subsection (c) is the clerical amendment; subsection (d) is the rule of construction;
and subsection (e) sets the effective date. Section 515D runs to subsections (a)
through (j), glossed next.}

\uscprov{1}{\textbf{Sec. 3, section 515D(a) Requirement.} The central rule. No
robot-patient interaction code may be generated or executed in a Physical AI oncology
clinical investigation unless an automated verification, validation, and uncertainty
quantification process, bound to one or more named external standards, has first
cleared for that interaction and the cleared record has been documented and attested
under subsection (e). This converts ``verify before generate'' from an engineering
best practice into a binding statutory precondition.}

\uscprov{1}{\textbf{Sec. 3, section 515D(b) Order of Operations.} Locks the sequence.
A sponsor may not generate, and may not execute, the code in advance of the cleared
verification record, and a record created afterward does not satisfy the section. This
forecloses generating and running code first and documenting later; the timing of the
record is itself the compliance fact.}

\uscprov{1}{\textbf{Sec. 3, section 515D(c) Gate Thresholds and Decision Rule.} The
quantitative test each interaction must pass. Every gate emits ACCEPT, BLOCK, or
ESCALATE: a single failing dimension sets BLOCK, and ambiguity or divergence sets
ESCALATE and returns the interaction to a qualified human. For each interaction the
verification fraction must equal 1.0, the validation agreement must meet or exceed the
per-gate threshold with the maximum relative error at or below the per-gate bound and
the human review recorded, the coefficient of variation across not fewer than three
seeded runs must stay within the per-gate bound, and any applicable hard catastrophe
predicate must also hold. The ten-gate schedule binds each gate to a minimum
agreement, a maximum coefficient of variation, the presence of a hard predicate, and a
governing external standard, making the gate auditable rather than discretionary;
gates 06, 09, and 10 carry the hard predicate at agreement 1.00 \cite{asmevv40,
iso14971}.}

\uscprov{1}{\textbf{Sec. 3, section 515D(d) Readiness Gates.} Site-level and
robot-level preconditions before a patient is enrolled or a robot acts. The trial site
must pass the Physical AI Standard Level gate on legislative authorization, regulatory
compliance, and operational readiness. Each robot must meet the minimum Unification
Standard Level for its functional type: not less than 7.0 for a surgical type, 6.0 for
a therapeutic positioning or diagnostic needle placement type, 4.0 for a rehabilitative
type, and 3.0 for a companion monitoring type. And each robot must be validated across
not fewer than two simulation frameworks within cross-framework tolerances of less than
2 millimeters of trajectory deviation and less than 0.5 newton of force discrepancy.
The graduated floor matches scrutiny to the physical risk of the function.}

\uscprov{1}{\textbf{Sec. 3, section 515D(e) Documentation and Attestation.} The record
subsection (a) requires. For each covered algorithm the sponsor files algorithm
documentation: a plain-language summary, the decision logic or architecture, the
training or conditioning data sets, the gate bindings identifying the governing
standard and threshold, the verification-before-generation record, and a determinism
statement giving the seed. A responsible officer signs a compliance attestation
certifying that the gate cleared before generation, that each gate met its
standard-bound threshold, that the site and robots cleared their readiness gates, that
bias on protected characteristics is minimized, and that a licensed human retained the
medical-necessity decision. All records sit in a hash-chained, tamper-evident audit
trail satisfying 21 CFR part 11, making the duty enforceable and the attestation a
personal certification \cite{cfr-part11}.}

\uscprov{1}{\textbf{Sec. 3, section 515D(f) Cybersecurity, Human Oversight, and
Lifecycle.} Operational controls across the system's life. Each Physical AI system
must incorporate authentication and access control, encryption in transit and at rest,
network segmentation, and software-integrity verification; operate under recorded human
oversight with a manual override and an emergency stop that reaches a safe state within
a procedure-appropriate budget plus hand-back-to-human escalation; and manage
configuration and model versions, monitor for model and data drift, and, on
decommissioning, archive the audit trail, securely erase patient data, and revoke
access. These duties guard the gate against tampering and degradation over time.}

\uscprov{1}{\textbf{Sec. 3, section 515D(g) Nondiscrimination.} Requires that each
covered algorithm and its training data minimize the risk of bias on protected
characteristics and adhere to evidence-based clinical guidelines, and that each system
be designed and operated consistent with applicable accessibility and accommodation
requirements. This carries Federal civil-rights and accessibility norms into the gate
as substantive criteria, not afterthoughts.}

\uscprov{1}{\textbf{Sec. 3, section 515D(h) Regulations.} Directs the Secretary to
issue final regulations not later than 365 days after enactment, and permits interim
guidance in the meantime, with a sponsor able to comply by meeting subsections (a)
through (g) until the regulations issue. This makes the statute self-executing rather
than dormant pending rulemaking.}

\uscprov{1}{\textbf{Sec. 3, section 515D(i) Definitions.} Defines the fifteen operative
terms, including Physical AI system, robot agent, robot-patient interaction code,
generative artificial intelligence, AI-enabled device software function, verification,
validation, uncertainty quantification, verification fraction (an equality to 1.0),
validation agreement, coefficient-of-variation bound (over not fewer than three seeded
runs), hard catastrophe predicate, Physical AI Standard Level, Unification Standard
Level (a 1.0 to 10.0 robot-readiness score), and hash-chained audit trail. Precise
definitions are what let the quantitative gate operate without discretion.}

\uscprov{1}{\textbf{Sec. 3, section 515D(j) Rule of Construction.} Confirms the
section's limits. It does not displace the investigational-device exemption under
section 520(g), human-subject protection under 21 CFR part 50, the IND framework under
21 CFR part 312, or the section 515C change-control authority as it applies outside
Physical AI, and it reclassifies no device. This positions section 515D as an additive
sequencing gate layered on the existing trial framework, not a replacement for it
\cite{cfr-part50, cfr-part312}.}

\uscprov{1}{\textbf{Sec. 3(b) Conforming Amendments (1) through (5).} Subsection (b)
threads the new duty through ten device provisions so the Act stays internally
consistent. Amendment (1), section 201(h) (21 U.S.C. 321(h)), confirms that software
that generates or executes robot-patient interaction code, or that directs the
autonomous physical action of a robot in contact with a human subject, is a device,
pulling the regulated code inside FDA jurisdiction. Amendment (2), section 301 (21
U.S.C. 331), adds a prohibited act reaching generation or execution of the code in
violation of section 515D and the failure to make or file the required documentation
and attestation, the enforcement hook. Amendment (3), section 501 (21 U.S.C. 351),
deems a Physical AI device software function used without a cleared section 515D record
adulterated. Amendment (4), section 505F (21 U.S.C. 355g), recognizes section 515D
records as real world evidence and applies the program to devices in Physical AI
oncology investigations. Amendment (5), section 510(k) (21 U.S.C. 360(k)), provides
that a verification-governed change consistent with an established predetermined change
control plan needs no new premarket notification \cite{usc-321h, usc-331, usc-351,
usc-355g, usc-360}.}

\uscprov{1}{\textbf{Sec. 3(b) Conforming Amendments (6) through (10).} Amendment (6),
section 513 (21 U.S.C. 360c), graduates special controls by degree of autonomy
(assistive, augmentative, or autonomous), matching rigor to how much the system acts on
its own. Amendment (7), section 515 (21 U.S.C. 360e), requires a premarket approval
application for such a device to include the cleared verification record and gate
evidence. Amendment (8), section 515C (21 U.S.C. 360e-4), requires a predetermined
change control plan for such a device to include an automated verification-before-change
protocol consistent with section 515D, carrying the rule forward into the lifecycle.
Amendment (9), section 520(o) (21 U.S.C. 360j(o)), confirms that such a software
function is not within the clinical-decision-support exclusion and remains a device,
closing the carve-out. Amendment (10), section 521 (21 U.S.C. 360k), adds a savings
clause preserving State human-in-the-loop and audit requirements that are not different
from, or in addition to, a Federal device requirement, protecting the State
human-over-AI laws cited in the findings \cite{usc-360c, usc-360e, usc-360e4, usc-360j,
usc-360k}.}

\uscprov{1}{\textbf{Sec. 3(c) Clerical Amendment.} Inserts the new ``Sec. 515D'' item
into the Act's table of contents after the item for section 515C, so the codified table
reflects the new section.}

\uscprov{1}{\textbf{Sec. 3(d) Rule of Construction.} Restates at the Act level that
nothing in the Act displaces the section 520(g) investigational-device exemption, 21
CFR part 50 human-subject protection, the 21 CFR part 312 IND framework, the section
515C authority outside Physical AI, or any State law preserved under section 521 as
amended, and that nothing reclassifies any device. It guards against reading the bill
more broadly than intended.}

\uscprov{1}{\textbf{Sec. 3(e) Effective Date; Implementation Deadline; Transition
Rule.} The Act takes effect on the first day of the first calendar quarter beginning
not less than one year after enactment; the Secretary must issue the final section
515D(h) regulations not later than 365 days after enactment; and an investigation
already enrolling on the effective date has 180 days to come into compliance, with the
February 2, 2026, effective date of the Quality Management System Regulation left
undisturbed. This gives sponsors and the agency a predictable runway \cite{fr-qmsr-2024}.}

\uscprov{1}{\textbf{Sec. 4. Comparative Print; Changes in Existing Law.} Section 4
reproduces the eleven affected Title 21 device sections in 21 U.S.C. order, showing
only the affected provisions, with matter to be deleted struck and matter to be
inserted marked. It is a Ramseyer-style comparative print, reproduced in full at
Appendix E: it carries no new duty but lets a reader see exactly how each conforming
amendment changes existing law, the transparency record a committee expects before
introduction.}
